\section{An exact finite lift for the reciprocity-fixed C2 sector}
\label{sec:c2-reciprocal-fixed-sector}

This section gives an evaluated affine semidefinite representation with
both inclusions proved.  Its scope is the full linear sector
$D=m$, $\alpha=C$ of the canonical C2 cone, including every zero
coefficient face.  It is not a converse for either of the previously
proposed two-variable Gram dictionaries outside this sector.

\paragraph{Complete parent and statement.}
For $c=(d,n,m,C,A,D,\alpha)\geq0$, put $D_u=1+u^2$, $D_v=1+v^2$ and
define $P_c=\psi^TM_c(u,v)\psi$, $\psi=(W,X,Y,Z)^T$, by
\begin{align*}
P_c={}&nD_uD_vX^2+m(vW-X+uvY-uZ)^2
 +CD_v(uW-Y)^2+AD_uD_vY^2\\
&+(D+d/2)(uW+uvX-Y-vZ)^2
 +(\alpha+d/2)D_u(vY-Z)^2\\
&+(d/2)D_v(X+uZ)^2-(d/2)D_uD_vZ^2.
\end{align*}
The canonical slack is $F_c=4P_c$.  Let $\mathcal C_2$ consist of these
nonnegative coefficient vectors for which $P_c\geq0$ for every real
angle and state, including homogeneous projective endpoints, and set
\[
\mathcal C_2^{\mathrm{fix}}
 =\mathcal C_2\cap\{D=m,\ \alpha=C\}.
\tag{SF1}
\]
The matrices $Q_+\in\mathbb S^4$, $Q_-\in\mathbb S^5$ in SF9--SF10
give the following exact representation:
\[
\boxed{\quad
c\in\mathcal C_2^{\mathrm{fix}}
\ \Longleftrightarrow\ 
\exists x\in\mathbb R^9:\
\operatorname{diag}\bigl(Q_+,Q_-,d,n,m,C,A,
 D-m,m-D,\alpha-C,C-\alpha\bigr)\succeq0.
\quad}
\tag{SF2}
\]
Thus the lift in the original seven coordinates has total matrix order
$18$ and nine auxiliary variables; the largest block has order $5$.
Using the five independent coordinates $(d,n,m,C,A)$ instead gives
total order $14$.  These are explicit sufficient sizes, with no claim
of minimality.

\paragraph{Why one angular curve is sufficient on this entire sector.}
For $U=u^2,V=v^2$, direct expansion gives
\[
\det M_c=D_u^2D_v^2UV L_c(U,V),\qquad
L_c=c_{11}+c_{21}U+c_{12}V+c_{20}U/V+c_{02}V/U+c_{10}/V+c_{01}/U.
\tag{SF3}
\]
On the sector $D=m,\alpha=C$, the coefficients satisfy
\[
\begin{aligned}
c_{21}=c_{10}=a_0
 &=\frac{AC(2m+d+2n)(2C+d+2m)}4,\\
c_{12}=c_{01}=b_0
 &=\frac{Cm(d+2n)(2A+2C+d+2m)}4.
\end{aligned}
\tag{SF4}
\]
In particular $a_0,b_0\geq0$, including on their zero faces.  Write
$U=rt$, $V=r/t$ with $r,t>0$.  All central and ratio terms in SF3
are independent of $r$, so the exact difference is
\[
L_c(rt,r/t)-L_c(t,1/t)
 =(r+r^{-1}-2)(a_0t+b_0/t)\geq0.
\tag{SF5}
\]
For $d>0$ and $uv\ne0$, the upper $WXY$ block of $M_c$ is positive
definite: the restriction of its signed core at $Z=0$ is
$\tfrac d2[(uW+uvX-Y)^2+D_uv^2Y^2+D_vX^2]$, whose three rows
are independent.  Its Schur complement therefore makes
$M_c(u,v)\succeq0$ equivalent to $L_c(U,V)\geq0$.
It follows from SF5 that testing $v=1/u$, $u>0$, for every state is
necessary and sufficient.  No state is restricted in taking this curve.
The exact sign gauges
\[
(u,v;W,X,Y,Z)\longmapsto
(\epsilon u,\eta v;W,\eta X,\epsilon Y,\epsilon\eta Z),
\qquad \epsilon,\eta\in\{-1,1\},
\]
preserve the complete parent.  They give both signs of the angles;
homogenization by the positive squares of the angle denominators and
continuity give all axes and projective endpoints.  If $d=0$, all six
remaining event packets are squares, so the same equivalence holds
directly.  Real symmetric positivity also gives complex-state
positivity by separating the real and imaginary parts of the state.

\paragraph{The exact univariate state polynomial.}
The matrix $H(u)=u^2M_c(u,1/u)$ is polynomial of degree four.
Set $W=S+T$ and $Z=S-T$, an invertible constant state transformation.
On SF1, the whole parent becomes
\[
u^2P_c(u,1/u;S+T,X,Y,S-T)
 =(S,X,Y)R(u)(S,X,Y)^T+q_T(u)T^2,
\tag{SF6}
\]
where
\[
q_T(u)=(1+u^2)\bigl[(2C+d+2m)u^2+2m\bigr]\geq0
\]
and the six independent entries of the symmetric matrix $R$ are
\begin{align*}
R_{SS}&=(2C+d+2m)u^4+(2C-d-4m)u^2+2m,\\
R_{SX}&=(d+2m)u^3-2mu,\\
R_{SY}&=-(2C+d+2m)u^3+(2m-2C)u,\\
R_{XX}&=nu^4+(2n+2m+d)u^2+n+d/2,\\
R_{XY}&=-(d/2+2m)u^2,\\
R_{YY}&=Au^4+(2A+2C+2m+d)u^2+A+2C+d/2.
\end{align*}
All state cross terms involving $T$ cancel identically, not only at a
stationary contact.  With $J_3=\operatorname{diag}(1,-1,-1)$,
\[
R(-u)=J_3R(u)J_3.
\tag{SF7}
\]
Consequently SF5--SF7 prove that membership in SF1 is equivalent to
$R(u)\succeq0$ for every real $u$.  At $u=0$ this follows by
continuity; no positive-definiteness assumption is imposed.

\paragraph{Finite degree and parity.}
The matrix Fej\'er--Riesz theorem on the real line states that a
Hermitian polynomial matrix of degree at most $2N$ which is positive
semidefinite on the real line is $G(u)^*G(u)$ with $\deg G\leq N$.
We use the degree-bounded version stated as Theorem 1.2 in
A. Zalar, \emph{Matrix Fej\'er--Riesz theorem with gaps},
arXiv:1503.06034v1, p.~1,
\url{https://arxiv.org/abs/1503.06034}.
Applied to the real symmetric $R$, splitting the real and imaginary
parts of $G$ gives real state-linear SOS rows of degree at most two
in $u$.  This also covers singular matrices.  A real-factor formulation
is given by C. Hanselka and R. Sinn, \emph{Positive Semidefinite
Univariate Matrix Polynomials}, arXiv:1707.08261,
\url{https://arxiv.org/abs/1707.08261}.

Average the resulting Gram identity under
$(u,S,X,Y)\mapsto(-u,S,-X,-Y)$, which preserves the target by SF7.
The nine degree-at-most-two state-linear monomials split into
\[
f_+=(S,uX,uY,u^2S)^T,\qquad
f_-=(X,Y,uS,u^2X,u^2Y)^T.
\tag{SF8}
\]
The average removes the off-diagonal block between the two parity
classes and preserves positive semidefiniteness.  Thus
$R(u)\succeq0$ for all real $u$ is equivalent to a Gram identity
\[
(S,X,Y)R(u)(S,X,Y)^T=f_+^TQ_+f_++f_-^TQ_-f_-,
\qquad Q_+\succeq0,\ Q_-\succeq0.
\]
Conversely such an identity is
plainly nonnegative for every state and real $u$.

\paragraph{Evaluated affine arrays.}
For arbitrary real $x_1,\ldots,x_9$, put
\[
Q_+=\begin{pmatrix}
2m&x_1&x_2&x_3\\
x_1&x_4&x_5&x_6\\
x_2&x_5&x_7&x_8\\
x_3&x_6&x_8&2C+d+2m
\end{pmatrix}.
\tag{SF9}
\]
Define the following eight explicitly affine entries:
\begin{align*}
\ell_1&=-2m-x_1,&\ell_2&=2m-2C-x_2,\\
\ell_3&=n+m+d/2-x_4/2,&\ell_4&=-d/2-2m-x_5-x_9,\\
\ell_5&=A+C+m+d/2-x_7/2,&\ell_6&=2C-d-4m-2x_3,\\
\ell_7&=d+2m-x_6,&\ell_8&=-2C-d-2m-x_8.
\end{align*}
Then
\[
Q_-=\begin{pmatrix}
n+d/2&0&\ell_1&\ell_3&x_9\\
0&A+2C+d/2&\ell_2&\ell_4&\ell_5\\
\ell_1&\ell_2&\ell_6&\ell_7&\ell_8\\
\ell_3&\ell_4&\ell_7&n&0\\
x_9&\ell_5&\ell_8&0&A
\end{pmatrix}.
\tag{SF10}
\]
For every choice of all five source variables and all nine auxiliaries,
these arrays satisfy the complete coefficient identity
\[
f_+^TQ_+f_++f_-^TQ_-f_-=(S,X,Y)R(u)(S,X,Y)^T.
\tag{SF11}
\]
They parametrize \emph{every} Gram identity in SF8.  Indeed the
$10+15=25$ symmetric Gram coordinates map onto sixteen independent
coefficients.  The coefficient map has rank sixteen; equivalently,
matching the powers $1,u,u^2,u^3,u^4$ in the six state pairs solves
every remaining entry in terms of precisely the nine entries chosen in
SF9--SF10.  The homogeneous auxiliary directions are independent,
so none of the allowed identities has been discarded.

\paragraph{Both inclusions and all boundary cases (SF12).}
If the pencil in SF2 is positive semidefinite, its scalar blocks force
the five independent coefficients to be nonnegative and force
$D=m$, $\alpha=C$.  SF11 gives $R(u)\succeq0$ for all real $u$;
SF6 gives positivity on $v=1/u$.  For $d>0$, SF3--SF5 and the
positive upper state block give positivity for all positive angle pairs,
and the gauges and homogeneous limits give every remaining test.
For $d=0$ the six event squares give the conclusion directly.
This proves the projected pencil is contained in $\mathcal C_2^{\rm fix}$.

Conversely a physical point in SF1 is positive on $v=1/u$ and every
state.  SF6 gives $R(u)\succeq0$ for $u>0$, and SF7 and continuity
extend this to all real $u$.  The degree-two matrix factorization and
parity average give positive semidefinite Gram blocks in SF8.
Coefficient-map completeness supplies nine real auxiliaries in exactly
SF9--SF10.  The orthant and equality scalar blocks also hold, proving
the other inclusion.  This argument includes zero matrices, vanishing
outer determinant masses, multiple zero event coefficients, and
singular Grams without taking inverses or strict limits.

\paragraph{Scope of the advance.}
The representation removes both angular quantifiers on the entire
reciprocity-fixed sector and evaluates all affine matrix coefficients.
It does not assert that SF8 is a two-angle SOS, nor that its Grams lie
in either the old or enlarged two-angle candidate.  For a general
source the equalities in SF4 fail, and the nonnegative projection
identity SF5 is unavailable.  The universal five-parameter regular
contact PSD-selection problem in GR17 remains to be proved.

The companion solver-free checker verifies the complete parent,
determinant identities, state congruence, parity, coefficient-map rank,
and every entry of the exported original-coordinate order-$18$ pencil.
The degree-bounded matrix factorization is the standard theorem quoted
above; it is not inferred from a solver result or from these finite checks.
